% K counts independent environments; n counts observations WITHIN each environment.
% Volume three architecture figures. Independent export only; not book fonts.
\colorlet{MutedBlueFill}{FigureBlueFill}
\colorlet{MutedGreenFill}{FigureGreenFill}
\colorlet{MutedPurple}{FigurePurple}
\colorlet{MutedPurpleFill}{FigurePurpleFill}
\tikzset{
 v3 picture/.style={x=1mm,y=1mm,font=\figurefont\fontsize{8}{10}\selectfont,text=FigureInk,line width=.65pt},
 v3 box/.style={draw=FigureBlue,fill=white,rounded corners=1.35mm,line width=.85pt,inner sep=2.5pt,align=center},
 v3 data/.style={v3 box,fill=FigureBlueFill},
 v3 model/.style={v3 box,draw=FigureGreen,fill=FigureGreenFill},
 v3 output/.style={v3 box,draw=FigureYellow,fill=FigureYellowFill},
 v3 risk/.style={v3 box,draw=FigureRed,fill=FigureRedFill},
 v3 group/.style={draw=FigureBlue,rounded corners=1.7mm,line width=.85pt,inner sep=3mm},
 v3 flow/.style={draw=FigureStroke,line width=.85pt,-{Stealth[length=1.9mm,width=1.25mm]},rounded corners=1.1mm},
 v3 guide/.style={draw=FigureMuted,line width=.6pt,dashed},
 v3 feedback/.style={v3 flow,draw=FigureGreen,dashed},
 v3 title/.style={draw=none,fill=none,font=\figurefont\fontsize{8}{10}\selectfont,align=center,inner sep=1pt},
 v3 note/.style={draw=none,fill=none,font=\FigureNoteFont\fontsize{7.5}{9.5}\selectfont,text=FigureMuted,align=center,inner sep=1pt},
 v3 label/.style={draw=none,fill=none,font=\figurefont\fontsize{8}{10}\selectfont,align=center,inner sep=1pt},
 v3 wire/.style={draw=FigureMuted,line width=.45pt},
 v3 point/.style={circle,draw=FigureBlue,fill=FigureBlueFill,minimum size=2.3mm,inner sep=0pt,line width=.55pt}
}
% Bring pre-existing native TikZ geometry to the same typography and scale.
\tikzset{
 foundation picture/.append style={font=\figurefont\fontsize{8}{10}\selectfont,line width=.65pt},
 foundation node/.append style={rounded corners=1.8pt,line width=.7pt,font=\figurefont\fontsize{8}{10}\selectfont},
 foundation heading/.append style={draw=none,fill=none,font=\figurefont\fontsize{8}{10}\selectfont},
 foundation note/.append style={draw=none,fill=none,font=\FigureNoteFont\fontsize{7.5}{9.5}\selectfont},
 foundation flow/.append style={line width=.75pt,-{Stealth[length=1.65mm,width=1.05mm]}},
 foundation frame/.append style={rounded corners=1.8pt,line width=.7pt},
 foundation edge label/.append style={draw=none,fill=none,font=\figurefont\fontsize{8}{10}\selectfont}
}
% Official Lucide outline icons, consistently used across diagrams. Size in mm.
\newcommand{\VIcon}[5]{\node[inner sep=0pt] at (#2,#3) {\includegraphics[width=#4mm]{figures/06-foundations/00-shared/tikz/icons/#1-#5.pdf}};}
% A neutral network glyph (not a claim about an actual model architecture).
% Coordinates denote its centre; scale 1 gives 13 x 8 mm.
\newcommand{\VNet}[4]{\begin{scope}[shift={(#1,#2)},scale=#3]
 \foreach \a in {-3,0,3}{\foreach \b in {-3,0,3}{\draw[v3 wire] (-5,\a)--(0,\b);}}
 \foreach \a in {-3,0,3}{\draw[v3 wire] (0,\a)--(5,0);}
 \foreach \a in {-3,0,3}{\draw[draw=#4,fill=#4,line width=.55pt] (-5,\a) circle(.9);\draw[draw=#4,fill=white,line width=.55pt] (0,\a) circle(.9);}
 \draw[draw=#4,fill=#4,line width=.55pt] (5,0) circle(.9);
\end{scope}}
\newcommand{\VDots}[4]{\foreach \k in {0,1,2,3}{\fill[#4] ({#1+2.7*\k},#2) circle(#3);}}
% Feature vector/tensor elements are scientific glyphs, not decorative icons.
\newcommand{\VVector}[4]{\foreach \k in {0,...,4}{\fill[#4,rounded corners=.25pt] ({#1+2.2*\k},#2) rectangle ++(1.5,#3);}}
% Detailed glyphs retain the visual vocabulary of the reference drawings.
% Layer counts are explicit: never substitute one topology for every network.
\colorlet{VNetworkInput}{FigureGreen}
\colorlet{VNetworkHidden}{FigureBlue}
\colorlet{VNetworkOutput}{FigureYellow}
\newcommand{\VLayerNetSpacing}[7]{\begin{scope}[shift={(#1,#2)},scale=#3]
 \foreach \i in {1,...,#4}{\foreach \j in {1,...,#5}{
  \draw[FigureStroke!85,line width=.68pt] (-7,{(\i-(#4+1)/2)*#7})--(0,{(\j-(#5+1)/2)*#7});}}
 \foreach \i in {1,...,#5}{\foreach \j in {1,...,#6}{
  \draw[FigureStroke!85,line width=.68pt] (0,{(\i-(#5+1)/2)*#7})--(7,{(\j-(#6+1)/2)*#7});}}
 \foreach \i in {1,...,#4}{\draw[draw=FigureStroke!85,fill=VNetworkInput!60!white,line width=.8pt] (-7,{(\i-(#4+1)/2)*#7}) circle(1.15);}
 \foreach \i in {1,...,#5}{\draw[draw=FigureStroke!85,fill=VNetworkHidden!48!white,line width=.8pt] (0,{(\i-(#5+1)/2)*#7}) circle(1.15);}
 \foreach \i in {1,...,#6}{\draw[draw=FigureStroke!85,fill=VNetworkOutput!70!white,line width=.8pt] (7,{(\i-(#6+1)/2)*#7}) circle(1.15);}
\end{scope}}
\newcommand{\VLayerNet}[6]{\VLayerNetSpacing{#1}{#2}{#3}{#4}{#5}{#6}{3.6}}
\newcommand{\VLayerNetTall}[6]{\VLayerNetSpacing{#1}{#2}{#3}{#4}{#5}{#6}{6.0}}
% Operators follow the DP-SGD visual grammar: a quiet rounded surround,
% a stronger circular core, and one readable mathematical operation.
% Arguments: node name, centre x/y, surround size (mm), semantic colour, math.
\newcommand{\VOperator}[6]{%
 \pgfmathsetmacro{\VOperatorCore}{#4*.74}%
 \node[draw=none,fill=#5!12!white,rounded corners=1.8pt,
   minimum width=#4mm,minimum height=#4mm,inner sep=0pt] (#1-outer) at (#2,#3) {};
 \node[circle,draw=#5,fill=#5!32!white,line width=.9pt,
   minimum size=\VOperatorCore mm,inner sep=.4pt,
   font=\figurefont\fontsize{10.5}{11.5}\selectfont,text=FigureInk] (#1) at (#2,#3) {$#6$};}
% Gaussian noise is a mechanism glyph, with the same surround and circle.
\newcommand{\VNoise}[5]{%
 \VOperator{#1}{#2}{#3}{#4}{#5}{}%
 \begin{scope}[shift={(#2,#3)},scale={#4/12}]
  \path[draw=#5,fill=#5!65!white,line width=.9pt]
    (-3,-1.9)--plot[domain=-3:3,samples=36] (\x,{-1.9+4.2*exp(-\x*\x/2)})--cycle;
 \end{scope}}
% Filled envelopes and folded sheets: thin strokes, distinct internal details.
\newcommand{\VEnvelope}[4]{\begin{scope}[shift={(#1,#2)},scale=#3,line join=round]
 \draw[draw=#4!85!FigureStroke,fill=#4!16!white,line width=.85pt,rounded corners=.7pt] (-6,-4) rectangle (6,4);
 \draw[draw=#4!85!FigureStroke,line width=.8pt] (-5.8,-3.7)--(-1,0.1) (5.8,-3.7)--(1,0.1);
 \draw[draw=#4!85!FigureStroke,fill=white,line width=.85pt] (-5.8,3.8)--(0,-.7)--(5.8,3.8);
\end{scope}}
\newcommand{\VDocument}[4]{\begin{scope}[shift={(#1,#2)},scale=#3,line join=round,line cap=round]
 \draw[draw=#4!85!FigureStroke,fill=#4!13!white,line width=.85pt,rounded corners=.4pt] (-4,-5.2)--(4,-5.2)--(4,2.5)--(1.5,5.2)--(-4,5.2)--cycle;
 \draw[draw=#4!85!FigureStroke,fill=#4!32!white,line width=.8pt] (1.5,5.2)--(1.5,2.5)--(4,2.5);
 \foreach \y/\len in {1.3/2.7,-.6/4.5,-2.5/4.5}{\draw[draw=#4!85!FigureStroke,line width=1pt] (-2.2,\y)--++(\len,0);}
\end{scope}}
\newcommand{\VDatabase}[4]{\begin{scope}[shift={(#1,#2)},scale=#3]
 \draw[draw=#4!85!FigureStroke,fill=#4!30!white,line width=.85pt] (-5.5,-4.5)--(-5.5,4.5) arc(180:360:5.5 and 1.8)--(5.5,-4.5) arc(0:-180:5.5 and 1.8);
 \draw[draw=#4!85!FigureStroke,line width=.8pt] (-5.5,-1.5) arc(180:360:5.5 and 1.8) (-5.5,1.5) arc(180:360:5.5 and 1.8);
 \draw[draw=#4!85!FigureStroke,fill=#4!48!white,line width=.85pt] (0,4.5) ellipse(5.5 and 1.8);
\end{scope}}
\newcommand{\VGear}[4]{\begin{scope}[shift={(#1,#2)},scale=#3]
 \path[fill=#4,rounded corners=.35pt] (-18:3.5)
  \foreach \a in {0,36,...,324}{--({\a-12}:3.5)--({\a-12}:4.4)--({\a+12}:4.4)--({\a+12}:3.5)--({\a+18}:3.5)}--cycle;
 \fill[white] (0,0) circle(1.65);
\end{scope}}
\newcommand{\VGears}[4]{\VGear{#1-2.3}{#2+1.1}{#3}{#4}\VGear{#1+4.2}{#2-3.0}{#3*.55}{#4}}
\newcommand{\VBrowser}[4]{\begin{scope}[shift={(#1,#2)},scale=#3,line cap=round]
 \draw[draw=#4!85!FigureStroke,fill=white,rounded corners=1pt,line width=.7pt] (-10,-7.5) rectangle (10,7.5);
 \path[fill=#4!35!white,rounded corners=1pt] (-9.7,4.6) rectangle (9.7,7.2);
 \draw[draw=#4!85!FigureStroke,line width=.65pt] (-10,4.5)--(10,4.5);
 \foreach \x in {-8.4,-6.8,-5.2}{\fill[#4!85!FigureStroke] (\x,5.9) circle(.4);}
 \fill[#4!12!white] (-9.7,-7.2) rectangle (-3.3,4.2);
 \foreach \y in {2.3,.5,-1.3}{\draw[#4!55!white,line width=1.6pt] (-8,\y)--(-5,\y);}
 \fill[#4!35!white,rounded corners=.5pt] (-1.8,-.2) rectangle (8.4,2.8);
 \fill[FigureGreen!35!white,rounded corners=.5pt] (4.7,-5.4) rectangle (8.4,-1.8);
 \foreach \y in {-2,-3.5,-5}{\draw[#4!40!white,line width=.9pt] (-1.8,\y)--(3.2,\y);}
\end{scope}}
\newcommand{\VRobot}[3]{\begin{scope}[shift={(#1,#2)},scale=#3,line join=round,line cap=round]
 \draw[draw=FigurePurple!85!FigureStroke,fill=FigurePurple!38!white,line width=.7pt] (-6,-7)--(-7,1)--(-1,6)--(4,2)--(2,0)--(-1,3)--(-4,0)--(-3,-7)--cycle;
 \foreach \x/\y/\r in {-5.5/0/1.8,-.7/5/1.65,3.8/1.5/1.25}{
  \draw[draw=FigurePurple!85!FigureStroke,fill=FigurePurple!45!white,line width=.7pt] (\x,\y) circle(\r);
  \fill[white] (\x,\y) circle(.65);}
 \draw[draw=FigureBlue!85!FigureStroke,fill=FigureBlue!35!white,line width=.65pt] (3.1,.4)--(2,-1.2)--(2.8,-2.9)--(3.5,-2.3)--(3,-1.3)--(4.2,-.2)--(5.2,-1.3)--(4.7,-2.3)--(5.4,-2.9)--(6.2,-1.2)--(5,.4);
 \draw[draw=FigurePurple!85!FigureStroke,fill=FigurePurple!38!white,rounded corners=.6pt,line width=.7pt] (-7.8,-8.6)--(-7.8,-7.1) .. controls (-7.8,-6) and (-1.2,-6) .. (-1.2,-7.1)--(-1.2,-8.6)--cycle;
\end{scope}}
\newcommand{\VPerson}[4]{\begin{scope}[shift={(#1,#2)},scale=#3]
 \fill[#4] (0,2.8) circle(1.9);
 \path[fill=#4,rounded corners=.5pt] (-3.3,-3.5)--(-3.3,-1.5) .. controls (-3.3,1.6) and (3.3,1.6) .. (3.3,-1.5)--(3.3,-3.5)--cycle;
\end{scope}}
\newcommand{\VServer}[4]{\begin{scope}[shift={(#1,#2)},scale=#3]
 \foreach \y in {-3,0,3}{
  \draw[draw=#4!85!FigureStroke,fill=#4!75!white,line width=.55pt,rounded corners=.6pt] (-5,\y-1.2) rectangle (5,\y+1.2);
  \fill[white] (3.8,\y) circle(.3);
  \draw[#4!30!white,line width=.55pt] (-3.5,\y)--(-.5,\y);}
\end{scope}}
\newcommand{\VMonitor}[4]{\begin{scope}[shift={(#1,#2)},scale=#3]
 \draw[draw=#4!85!FigureStroke,fill=#4!25!white,line width=.7pt,rounded corners=.6pt] (-6,-3.6) rectangle (6,4.3);
 \fill[white] (-5.2,-2.7) rectangle (5.2,3.5);
 \fill[#4!85!FigureStroke,rounded corners=.4pt] (-.8,-5.3) rectangle (.8,-3.5);
 \draw[#4!85!FigureStroke,line width=1pt,line cap=round] (-3.3,-5.3)--(3.3,-5.3);
\end{scope}}

\begin{tikzpicture}[v3 picture]
\tikzset{v3 box/.append style={line width=.8pt},v3 group/.append style={line width=.85pt},v3 flow/.append style={line width=.9pt}}
\path[use as bounding box] (0,0) rectangle (163,94);
\node[v3 title] at (55,88) {原始样本};
\node[v3 title] at (97,88) {增加环境内样本};
\node[v3 title,text=FigureYellow] at (142,88) {增加独立环境};
% Restore three comparison groups; headings sit outside, with no title band.
\draw[v3 group,fill=FigureBlue!3!white,line width=.85pt] (38.5,42.5) rectangle (75.5,82);
\draw[v3 group,fill=FigureBlue!3!white,line width=.85pt] (79.5,42.5) rectangle (116.5,82);
\draw[v3 group,draw=FigureYellow,fill=FigureYellow!4!white,line width=.85pt] (124.5,3.5) rectangle (161.5,82);
% Distribution of environments. Its graph does not assert causal relationships.
\node[v3 label] at (14,65) {环境分布 $\mathcal Q$};
\path[draw=FigureBlue,fill=FigureBlue!8!white,line width=1.1pt,rounded corners=.6pt]
 (3,46) .. controls (0,49) and (2,53) .. (5,54)
 .. controls (5,58) and (8,60) .. (11,59)
 .. controls (14,63) and (18,62) .. (20,58)
 .. controls (25,58) and (28,56) .. (27,52)
 .. controls (30,48) and (27,43) .. (22,43)
 .. controls (20,41) and (19,40.5) .. (18,41)
 .. controls (16,39.5) and (13,40) .. (12,42)
 .. controls (7,40) and (3,42) .. (3,46)--cycle;
% Draw the undirected environment graph first, then every node: no edge paints across a node interior.
\foreach \x/\y in {8/49,14/55,10/44,21/50,20/44}{\draw[v3 wire,line width=.8pt] (14,48)--(\x,\y);}
\foreach \x/\y in {8/49,14/55,14/48,10/44,21/50,20/44}{\draw[FigureBlue,fill=FigureBlue!60!white,line width=.8pt] (\x,\y) circle(1.3);}
\foreach \x/\n in {40/3,81/6,126/3}{
 \foreach \r/\y in {1/75,2/62,3/49}{
  \draw[v3 box,fill=FigureBlue!8!white] (\x,{\y-5}) rectangle ++(34,10);
  \node[v3 label,anchor=west] at ({\x+4},\y) {$P_{\r}$};
  \ifnum\n=3\relax
   \foreach \d in {11,20,29}{\fill[FigureBlue] ({\x+\d},\y) circle(1.5);}
  \else
   \foreach \d in {10,14,18}{\fill[FigureBlue] ({\x+\d},\y) circle(1.2);}
   \foreach \d in {22,26,30}{\draw[FigureBlue,fill=white,line width=.8pt] ({\x+\d},\y) circle(1.2);}
  \fi
 }
}
% Blue environments P1,P2,P3 are the SAME in all three columns.
% One cloud-boundary outlet forks into the three original environment samples.
\draw[draw=FigureStroke,line width=.95pt] (27,52)--(32,52);
\draw[draw=FigureStroke,line width=.95pt] (32,49)--(32,75);
\foreach \y in {75,62,49}{\draw[v3 flow,line width=.95pt] (32,\y)--(40,\y);}
\foreach \y in {52,62}{\fill[FigureStroke] (32,\y) circle(.4);}
\foreach \y in {75,62,49}{\draw[v3 flow,line width=.95pt,rounded corners=2mm] (74,\y)--(81,\y);}
\node[v3 label] at (57,38) {$K=3,\ n=3$};\node[v3 note] at (57,32) {$9$ 个观测};
\node[v3 label] at (98,38) {$K=3,\ n=6$};\node[v3 note] at (98,32) {$18$ 个观测};
% The common trunk represents three independent draws, not three copies of one P.
\foreach \r/\y in {4/36,5/23,6/10}{
 \draw[v3 box,draw=FigureGreen,fill=FigureGreen!8!white] (126,{\y-5}) rectangle ++(34,10);
 \node[v3 label,anchor=west] at (130,\y) {$P_{\r}$};
 \foreach \d in {137,146,155}{\fill[FigureBlue] (\d,\y) circle(1.5);}
}
\draw[draw=FigureStroke,line width=.95pt,rounded corners=2mm] (18,41)--(18,23)--(121,23);
\draw[draw=FigureStroke,line width=.95pt] (121,10)--(121,36);
\fill[FigureStroke] (121,23) circle(.4);
\foreach \y in {36,23,10}{\draw[v3 flow,line width=.95pt] (121,\y)--(126,\y);}
\node[v3 note,anchor=south] at (73,25) {$P_4,P_5,P_6\sim\mathcal Q$，独立抽取};
\node[v3 label] at (142,1) {$K=6,\ n=3$：$18$ 个观测};
\end{tikzpicture}
